\documentclass[12pt,a4paper]{article}

% =====================================================================
% PACOTES
% =====================================================================
\usepackage[brazilian]{babel}
\usepackage[T1]{fontenc}
\usepackage{amsmath}
\usepackage{amssymb}
\usepackage{amsthm}
\usepackage{geometry}
\usepackage{booktabs}
\usepackage{listings}
\usepackage{xcolor}
\usepackage{hyperref}
\usepackage{mathtools}
\usepackage{enumitem}

\geometry{top=2.5cm, bottom=2.5cm, left=3cm, right=2.5cm}

% =====================================================================
% CONFIGURACAO DO LISTINGS (codigo C)
% =====================================================================
\lstdefinestyle{cstyle}{
  language=C,
  basicstyle=\ttfamily\footnotesize,
  keywordstyle=\color{blue!80!black}\bfseries,
  commentstyle=\color{green!50!black}\itshape,
  stringstyle=\color{red!60!black},
  numbers=left,
  numberstyle=\tiny\color{gray},
  stepnumber=1,
  numbersep=10pt,
  frame=single,
  rulecolor=\color{black!25},
  breaklines=true,
  tabsize=4,
  showstringspaces=false,
  captionpos=b,
}
\lstset{style=cstyle}

% =====================================================================
% AMBIENTES
% =====================================================================
\theoremstyle{definition}
\newtheorem{questao}{Quest\~ao}

\newenvironment{solucao}{%
  \par\medskip\noindent\textbf{Solu\c{c}\~ao.}\quad\ignorespaces
}{%
  \par\medskip
}

\newenvironment{critica}{%
  \par\medskip\noindent\textbf{Cr\'itica.}\quad\ignorespaces
}{%
  \par\medskip
}

% caixa de resultado
\newcommand{\resultado}[1]{%
  \begin{center}
    \fbox{\parbox{0.85\linewidth}{\centering $\displaystyle #1$}}
  \end{center}
}

% =====================================================================
% HYPERREF
% =====================================================================
\hypersetup{
  colorlinks=true,
  linkcolor=blue!70!black,
  urlcolor=blue,
  pdftitle={Resolucao Prova 1 -- PMMC CEFET-MG 2013},
}

% =====================================================================
% CABECALHO
% =====================================================================
\title{%
  \textbf{Resolu\c{c}\~ao Comentada -- Prova 1}\\[0.6ex]
  \large Princ\'ipios de Modelagem Matem\'atica e Computacional\\[0.3ex]
  \normalsize CEFET-MG $\cdot$ Programa de P\'os-Gradua\c{c}\~ao em
  Modelagem Matem\'atica e Computacional\\
  II\textsuperscript{o} semestre de 2025 $\cdot$ Professor Allbens Atman
}
\author{}
\date{\today}

% =====================================================================
\begin{document}
% =====================================================================

\maketitle
\tableofcontents
\newpage

% ===================================================================
% QUESTAO 1
% ===================================================================
\section{Quest\~ao 1 -- Dimens\~ao Fractal}

\subsection{F\'ormula geral de auto-similaridade}

Para um fractal \textbf{auto-similar}, a dimens\~ao de Hausdorff--Besicovitch \'e
definida como o limite:
\begin{equation}
  D_f = \lim_{k\to\infty} \frac{\ln N_p^k}{\ln(1/p)^k}
      = \frac{\ln N_p}{\ln(1/p)},
  \label{eq:df}
\end{equation}
onde:
\begin{itemize}
  \item $N_p$ = n\'umero de c\'opias auto-similares em cada itera\c{c}\~ao;
  \item $p$ = fator de escala linear (comprimento da c\'opia / comprimento original);
  \item $1/p$ = raz\~ao de auto-similaridade.
\end{itemize}

% ------- Conjunto A -------
\subsection{Conjunto~A -- Curva tipo Koch}

A regra geradora da curva mostrada no enunciado \'e:
\begin{quote}
  Cada segmento de comprimento $\ell$ \'e substitu\'ido por
  \textbf{4 c\'opias} de comprimento $\ell/3$.
\end{quote}
Logo: $N_p = 4$, $p = 1/3$.

\paragraph{Verifica\c{c}\~ao iterativa.}

\begin{center}
\begin{tabular}{ccc}
\toprule
Itera\c{c}\~ao $k$ & N\'umero de segmentos & Comprimento total\\
\midrule
0 & 1             & $\ell$ \\
1 & 4             & $(4/3)\,\ell$ \\
2 & 16            & $(4/3)^2\,\ell$ \\
$k$ & $4^k$       & $(4/3)^k\,\ell$ \\
\bottomrule
\end{tabular}
\end{center}

O comprimento total cresce sem limite ($\to\infty$), confirmando dimensionalidade
n\~ao inteira.

\paragraph{C\'alculo.}
\begin{equation}
  D_f = \frac{\ln 4}{\ln 3}
      = \frac{2\ln 2}{\ln 3}
      \approx \frac{1{,}386}{1{,}099}
\end{equation}
\resultado{D_f \approx 1{,}26}

\paragraph{Interpreta\c{c}\~ao.}
$1 < D_f < 2$: o conjunto \'e mais complexo que uma curva ($D=1$), mas n\~ao
preenche o plano ($D=2$).
O per\'imetro diverge, mas a medida de \'area \'e nula.

% ------- Conjunto B -------
\subsection{Conjunto~B -- Pir\^amide de Sierpinski (3D)}

A sequ\^encia mostrada \'e o \textbf{tetraedro (pir\^amide) de Sierpinski}:
\begin{quote}
  Cada tetraedro \'e substitu\'ido por \textbf{4 c\'opias} em escala $1/2$.
\end{quote}
Logo: $N_p = 4$, $p = 1/2$.

\paragraph{Verifica\c{c}\~ao iterativa.}
\begin{center}
\begin{tabular}{ccc}
\toprule
Itera\c{c}\~ao $k$ & N\'umero de c\'opias $N$ & Fator de escala\\
\midrule
1 & $4^1 = 4$    & $1/2$ \\
2 & $4^2 = 16$   & $1/4$ \\
$k$ & $4^k$      & $1/2^k$ \\
\bottomrule
\end{tabular}
\end{center}

\paragraph{C\'alculo.}
\begin{equation}
  D_f = \frac{\ln 4}{\ln 2}
      = \frac{2\ln 2}{\ln 2} = 2.
\end{equation}
\resultado{D_f = 2}

\paragraph{Interpreta\c{c}\~ao.}
$D_f = 2$: o tetraedro de Sierpinski tem dimens\~ao de \emph{superf\'icie},
igual \`a de um plano s\'olido, apesar de ser um objeto tridimensional poroso.
Isso ocorre porque o n\'umero de c\'opias ($N=4$) exatamente compensa a redu\c{c}\~ao
de escala ($p = 1/2$ em 3D implicaria $N = 8$ para preencher o volume;
com $N=4$ o volume \'e nulo e apenas a superf\'icie persiste).

% ===================================================================
% QUESTAO 2
% ===================================================================
\section{Quest\~ao 2 -- Grua no Est\'adio: Frequ\^encia de Oscila\c{c}\~ao}

\subsection{Identifica\c{c}\~ao das vari\'aveis e dimens\~oes}

As grandezas envolvidas e suas dimens\~oes em (M, L, T) s\~ao:

\begin{center}
\begin{tabular}{lll}
\toprule
\textbf{S\'imbolo} & \textbf{Grandeza} & \textbf{Dimens\~ao}\\
\midrule
$\omega$   & frequ\^encia de oscila\c{c}\~ao (incognita) & $T^{-1}$ \\
$m$        & massa da viga                                 & $M$ \\
$g$        & acelera\c{c}\~ao gravitacional                & $LT^{-2}$ \\
$\tau$     & tens\~ao no cabo (for\c{c}a)                  & $MLT^{-2}$ \\
$L$        & comprimento do cabo                           & $L$ \\
$h$        & altura da viga acima do solo                  & $L$ \\
\bottomrule
\end{tabular}
\end{center}

\noindent
\textit{Nota}: ``tens\~ao no cabo'' \'e interpretada como \textbf{for\c{c}a} ($[\tau] = MLT^{-2}$),
n\~ao como tens\~ao mec\^anica (for\c{c}a/\'area).

\subsection{Contagem Pi de Buckingham}

\begin{align}
  n &= 6 \quad \text{(vari\'aveis)}, \\
  r &= 3 \quad \text{(dimens\~oes fundamentais: M, L, T)}, \\
  n - r &= 3 \quad \text{(grupos adimensionais $\Pi_i$)}.
\end{align}

\paragraph{Vari\'aveis repetidoras.}
Escolhemos $\{\omega,\, h,\, m\}$ como repetidoras, pois cobrem T, L e M
de forma independente:
$[\omega] = T^{-1}$, $[h] = L$, $[m] = M$.

\subsection{Constru\c{c}\~ao dos grupos \texorpdfstring{$\Pi_i$}{Pi\_i}}

Cada grupo tem a forma $\Pi_i = \omega^{a_i}\,h^{b_i}\,m^{c_i}\,Q_i$,
onde $Q_i$ \'e a vari\'avel extra.

\paragraph{Grupo $\Pi_1$ com $Q_1 = L$ (comprimento do cabo).}
\begin{equation}
  \Pi_1 = \omega^{a_1}\,h^{b_1}\,m^{c_1}\,L
  \implies T^{-a_1}\,L^{b_1+1}\,M^{c_1} = M^0 L^0 T^0.
\end{equation}
Sistema:
$\;c_1 = 0,\quad b_1 + 1 = 0 \Rightarrow b_1 = -1,\quad a_1 = 0.$
\resultado{\Pi_1 = \frac{L}{h}}

\paragraph{Grupo $\Pi_2$ com $Q_2 = g$.}
\begin{equation}
  \Pi_2 = \omega^{a_2}\,h^{b_2}\,m^{c_2}\,g
  \implies M^{c_2}\,L^{b_2+1}\,T^{-a_2-2} = 1.
\end{equation}
Sistema:
$\;c_2 = 0,\quad b_2 = -1,\quad a_2 = -2.$
\resultado{\Pi_2 = \frac{g}{h\,\omega^2}}

\paragraph{Grupo $\Pi_3$ com $Q_3 = \tau$ ($[\tau]=MLT^{-2}$).}
\begin{equation}
  \Pi_3 = \omega^{a_3}\,h^{b_3}\,m^{c_3}\,\tau
  \implies M^{c_3+1}\,L^{b_3+1}\,T^{-a_3-2} = 1.
\end{equation}
Sistema:
$\;c_3 = -1,\quad b_3 = -1,\quad a_3 = -2.$
\resultado{\Pi_3 = \frac{\tau}{m\,h\,\omega^2}}

\subsection{Resultado final}

O Teorema Pi garante $f(\Pi_1,\,\Pi_2,\,\Pi_3) = 0$.
Das rela\c{c}\~oes inversas:
\begin{align}
  \Pi_2 = \text{const.}
  &\implies \boxed{\omega \propto \sqrt{\frac{g}{h}}},
  \label{eq:omega1}\\[6pt]
  \Pi_3 = \text{const.}
  &\implies \boxed{\omega \propto \sqrt{\frac{\tau}{m\,h}}}.
  \label{eq:omega2}
\end{align}

\paragraph{Interpreta\c{c}\~ao f\'isica.}
\begin{itemize}
  \item A rela\c{c}\~ao \eqref{eq:omega1} \'e \textbf{an\'aloga ao p\^endulo simples}
    ($T = 2\pi\sqrt{h/g}$): a viga suspensa a altura $h$ oscila como um p\^endulo
    de comprimento $h$. Quanto maior a altura, menor a frequ\^encia.
  \item A rela\c{c}\~ao \eqref{eq:omega2} mostra que maior tens\~ao $\tau$
    (cabo mais retesado) ou menor massa $m$ aumentam $\omega$ --- coerente com
    vibrações de cordas.
  \item O grupo $\Pi_1 = L/h$ \'e adimensional geom\'etrico: a raz\~ao entre
    comprimento do cabo e altura influencia o modo de vibra\c{c}\~ao,
    mas n\~ao aparece diretamente na estimativa de ordem de grandeza.
\end{itemize}

% ===================================================================
% QUESTAO 3
% ===================================================================
\section{Quest\~ao 3 -- Deprecia\c{c}\~ao do Carro (Fun\c{c}\~ao Exponencial)}

\textbf{Dados:} valor inicial $V_0 = \text{R\$}\;51{.}500$; perda de $12\%$ do valor a cada ano.

\subsection{a) Fun\c{c}\~ao do valor ap\'os \texorpdfstring{$t$}{t} anos}

A cada ano o carro ret\'em $(1 - 0{,}12) = 0{,}88$ do valor anterior.
Aplicando $t$ vezes:
\begin{equation}
  \boxed{V(t) = 51{.}500 \times (0{,}88)^t}, \qquad t \ge 0.
  \label{eq:carro}
\end{equation}

\paragraph{Tabela de valores.}
\begin{center}
\begin{tabular}{rr}
\toprule
$t$ (anos) & $V(t)$ (R\$) \\
\midrule
0 & 51.500,00 \\
1 & 45.320,00 \\
2 & 39.881,60 \\
3 & 35.096,21 \\
4 & 30.884,66 \\
6 & 23.916,81 \\
\bottomrule
\end{tabular}
\end{center}

\subsection{b) Taxa percentual mensal aproximada}

Um m\^es corresponde a $t = 1/12$ de ano.
Chamando $i_m$ a taxa mensal de deprecia\c{c}\~ao:
\begin{equation}
  (1 - i_m)^{12} = 0{,}88
  \implies
  1 - i_m = 0{,}88^{1/12}.
\end{equation}

C\'alculo exato:
\begin{equation}
  0{,}88^{1/12} = e^{\frac{\ln 0{,}88}{12}} = e^{-0{,}12783/12} \approx e^{-0{,}010653} \approx 0{,}98941.
\end{equation}
\resultado{i_m = 1 - 0{,}98941 \approx 1{,}059\% \approx 1{,}06\%\text{ ao m\^es}}

\paragraph{Aproxima\c{c}\~ao binomial (primeira ordem).}
Para $\ln(1-x)\approx -x$ com $x = 0{,}12$ pequeno:
\begin{equation}
  i_m \approx \frac{0{,}12}{12} = 1\%\text{ ao m\^es}.
\end{equation}
A diferen\c{c}a em rela\c{c}\~ao ao valor exato ($1{,}06\%$) \'e inferior a $6\%$ relativo ---
aceit\'avel como estimativa de ordem de grandeza.

\subsection{c) Gr\'afico da fun\c{c}\~ao}

O gr\'afico de $V(t) = 51{.}500 \times 0{,}88^t$ \'e uma \textbf{exponencial decrescente}:
assintota horizontal em $V = 0$, valor inicial $V(0) = 51{.}500$ e
concavidade positiva em todo o dom\'inio.

Em escala \textbf{semi-logar\'itmica} $(\log_{10} V$ \textit{vs} $t)$:
\begin{equation}
  \log_{10} V = \log_{10} 51{.}500 + t\,\log_{10}(0{,}88)
              \approx 4{,}7118 - 0{,}0555\,t.
\end{equation}
O gr\'afico torna-se uma \textbf{reta de inclinac\c{c}\~ao} $\approx -0{,}056\,\text{dec}^{-1}$.

\begin{center}
\begin{tabular}{rll}
\toprule
$t$ & $V(t)$ (R\$) & $\log_{10} V$ \\
\midrule
0 & 51.500 & 4,712 \\
1 & 45.320 & 4,656 \\
2 & 39.882 & 4,601 \\
3 & 35.096 & 4,545 \\
4 & 30.885 & 4,490 \\
6 & 23.917 & 4,379 \\
\bottomrule
\end{tabular}
\end{center}

\subsection{d) Estimativa do valor ap\'os 6 anos}

Diretamente da equa\c{c}\~ao \eqref{eq:carro}:
\begin{equation}
  V(6) = 51{.}500 \times (0{,}88)^6
        = 51{.}500 \times 0{,}46434
        \approx \text{R\$ } 23{.}916{,}81.
\end{equation}

Pelo gr\'afico semi-log: para $t = 6$, $\log_{10} V \approx 4{,}35$,
logo $V \approx 10^{4{,}35} \approx \text{R\$}\;22{.}387$.
A pequena diferen\c{c}a deve-se \`a interpola\c{c}\~ao gr\'afica.

\resultado{V(6) \approx \text{R\$ }23{.}917}

% ===================================================================
% QUESTAO 4
% ===================================================================
\section{Quest\~ao 4 -- Gerador Linear Congruencial (N\'umeros M\'agicos)}

\subsection{Defini\c{c}\~ao e Teorema de Hull-Dobell}

O gerador linear congruencial (LCG) \'e definido pela recorr\^encia:
\begin{equation}
  X_{n+1} = (a\,X_n + c)\;\bmod\;m,
  \label{eq:lcg}
\end{equation}
onde $m$ \'e o m\'odulo, $a$ o multiplicador e $c$ o incremento.

\paragraph{Condi\c{c}\~oes para per\'iodo m\'aximo $m$ (Hull-Dobell, 1962).}
\begin{enumerate}
  \item $\gcd(c,\,m) = 1$ \quad ($c$ e $m$ coprimos);
  \item $a - 1$ divisível por todos os fatores primos de $m$;
  \item Se $m$ é múltiplo de $4$, então $a-1$ também deve ser múltiplo de $4$.
\end{enumerate}

\subsection{a) Par m\'agico para \texorpdfstring{$m = 10$}{m = 10}}

Fatoração: $m = 10 = 2 \times 5$ (fatores primos: $2$ e $5$; $m$ \textbf{n\~ao} \'e m\'ultiplo de $4$,
portanto a condição 3 n\~ao se aplica).

\paragraph{Escolha de $a$.}
A condição 2 exige que $a - 1$ seja divis\'ivel por $2$ e por $5$,
ou seja, $a - 1$ dev\'e ser m\'ultiplo de $10$:
\begin{equation}
  a \in \{11, 21, 31, 41, \ldots\}.
\end{equation}

\paragraph{Escolha de $c$.}
A condição 1 exige $\gcd(c, 10) = 1$, ou seja, $c$ n\~ao pode ser m\'ultiplo de $2$ nem de $5$:
\begin{equation}
  c \in \{1, 3, 7, 9, 11, \ldots\}.
\end{equation}

\paragraph{Par v\'alido.}
Tomando $a = 31$ e $c = 9$ (como na folha da prova):
\begin{itemize}
  \item $a - 1 = 30 = 2 \times 3 \times 5$ --- divis\'ivel por $2$ e $5$ \checkmark
  \item $\gcd(9,\,10) = 1$ \checkmark
  \item $m = 10$ n\~ao \'e m\'ultiplo de $4$ (condição 3 n\~ao se aplica) \checkmark
\end{itemize}
\resultado{(a,c) = (31,\;9)}

O gerador fica:
\begin{equation}
  X_{n+1} = (31\,X_n + 9)\;\bmod\;10.
\end{equation}

\subsection{b) Demonstra\c{c}\~ao com tr\^es sementes}

\paragraph{Nota de c\'alculo.}
Como $31 \equiv 1\pmod{10}$, tem-se $(31\,X_n + 9)\bmod 10 = (X_n + 9)\bmod 10$.
O incremento $+9 \equiv -1\pmod{10}$, logo cada passo \textbf{decrementa} o d\'igito em $1$.

\paragraph{Semente $X_0 = 0$.}
\begin{equation}
  0 \to 9 \to 8 \to 7 \to 6 \to 5 \to 4 \to 3 \to 2 \to 1 \to 0 \quad
  (\text{per\'iodo } = 10\ \checkmark)
\end{equation}

\paragraph{Semente $X_0 = 3$.}
\begin{equation}
  3 \to 2 \to 1 \to 0 \to 9 \to 8 \to 7 \to 6 \to 5 \to 4 \to 3 \quad
  (\text{per\'iodo } = 10\ \checkmark)
\end{equation}

\paragraph{Semente $X_0 = 7$.}
\begin{equation}
  7 \to 6 \to 5 \to 4 \to 3 \to 2 \to 1 \to 0 \to 9 \to 8 \to 7 \quad
  (\text{per\'iodo } = 10\ \checkmark)
\end{equation}

Em todos os casos o gerador percorre os $10$ res\'iduos $\{0,1,\ldots,9\}$ antes de repetir,
confirmando per\'iodo m\'aximo $m = 10$ para qualquer semente inteira.

\begin{center}
\begin{tabular}{cccc}
\toprule
Passo $n$ & $X_n$ (semente 0) & $X_n$ (semente 3) & $X_n$ (semente 7) \\
\midrule
0 & 0 & 3 & 7 \\
1 & 9 & 2 & 6 \\
2 & 8 & 1 & 5 \\
3 & 7 & 0 & 4 \\
4 & 6 & 9 & 3 \\
5 & 5 & 8 & 2 \\
6 & 4 & 7 & 1 \\
7 & 3 & 6 & 0 \\
8 & 2 & 5 & 9 \\
9 & 1 & 4 & 8 \\
10 & 0 & 3 & 7 \\
\bottomrule
\end{tabular}
\end{center}

% ===================================================================
% TOPICOS DE ESTUDO COMPLEMENTAR
% ===================================================================
\newpage
\section{T\'opicos de Estudo Complementar}

Esta se\c{c}\~ao re\'une os t\'opicos citados para estudo que n\~ao compuseram diretamente
as quatro quest\~oes da prova, mas s\~ao recorrentes nas avalia\c{c}\~oes da disciplina.

% --- 5.1 ---
\subsection{Deca\-imento por Carbono-14}

\textbf{Enunciado (Prova -- II\textordmasculine{} semestre 2013).}
Uma amostra \'e datada em $570\,000$ anos. Concentra\c{c}\~ao natural: $1\,\text{\'atomo em }10^{12}$;
meia-vida do ${}^{14}$C: $t_{1/2} = 5700$ anos.
\begin{enumerate}[label=\alph*)]
  \item Concentra\c{c}\~ao de ${}^{14}$C na amostra.
  \item Incerteza na idade se a precis\~ao \'e $\delta C/C = 10^{-23}$.
  \item Cr\'itica ao resultado.
\end{enumerate}

\paragraph{Modelo.}
\begin{equation}
  C(t) = C_0\,e^{-\lambda t}, \qquad \lambda = \frac{\ln 2}{t_{1/2}}
         = \frac{\ln 2}{5700}\;\text{anos}^{-1}.
\end{equation}

\paragraph{a) Concentra\c{c}\~ao ap\'os 570\,000 anos.}
\begin{equation}
  n_{1/2} = \frac{570\,000}{5\,700} = 100 \text{ meias-vidas}.
\end{equation}
Como $2^{10}\approx10^{3}$, temos $2^{100}\approx10^{30}$, logo:
\begin{equation}
  C = \frac{1}{10^{12}}\cdot\frac{1}{2^{100}} \approx \frac{1}{10^{42}}\;\text{\'atomos por \'atomo de carbono.}
\end{equation}
\resultado{C \approx 10^{-42}}

\paragraph{b) Incerteza na idade.}
De $t = -\ln(C/C_0)/\lambda$:
\begin{equation}
  \delta t = \frac{t_{1/2}}{\ln 2}\,\frac{\delta C}{C} \approx 8225 \times 10^{-23}
           \approx 8\times10^{-20}\;\text{anos.}
\end{equation}
\resultado{\delta t \approx 8\times10^{-20}\,\text{anos (infinitesimal)}}

\paragraph{c) Cr\'itica.}
A concentra\c{c}\~ao esperada ($\sim 10^{-42}$) \'e 19 ordens de grandeza abaixo da
precis\~ao instrumental ($\sim 10^{-23}$).
O m\'etodo ${}^{14}$C \'e confi\'avel somente at\'e $\approx 50\,000$ anos ($\sim$8--9 meias-vidas).
Qualquer sinal detectado em amostras mais antigas \'e contaminação moderna, tornando
a datação de 570\,000 anos f\'isicamente imposs\'ivel por este m\'etodo.

% --- 5.2 ---
\subsection{Lista 4.10 -- Peso a Diferentes Altitudes}

\textbf{Enunciado.} Um corpo pesa $W_0 = 10\,\text{N}$ na superf\'icie da Terra.
Qual \'e seu peso a $h = 10\,\text{m}$ e no pico do Monte Everest ($h = 8\,848\,\text{m}$)?

\paragraph{Modelo (Newton).}
\begin{equation}
  W(h) = W_0\!\left(\frac{R}{R+h}\right)^{2}, \qquad R = 6{,}371\times10^{6}\,\text{m}.
\end{equation}

\paragraph{Resultado.}
\begin{center}
\begin{tabular}{lll}
\toprule
Altitude & Peso & Varia\c{c}\~ao\\
\midrule
Superf\'icie  & $10{,}0000\,\text{N}$ & ---\\
$h=10\,\text{m}$ & $9{,}99997\,\text{N}$ & $-0{,}0003\,\%$\\
Everest ($h=8848\,\text{m}$) & $9{,}9723\,\text{N}$ & $-0{,}277\,\%$\\
\bottomrule
\end{tabular}
\end{center}

\paragraph{Liga\c{c}\~ao com Taylor.}
Para $h\ll R$, usando $(1+x)^{-2}\approx 1-2x$:
\begin{equation}
  W(h)\approx W_0\!\left(1 - \frac{2h}{R}\right),
\end{equation}
reproduzindo a varia\c{c}\~ao linear com a altitude de primeira ordem.

% --- 5.3 ---
\subsection{Lista 4.24 -- Varia\c{c}\~ao do Per\'iodo do P\^endulo}

\textbf{Enunciado.} Em que porcentagem o per\'iodo de um p\^endulo muda se seu
comprimento for reduzido para $L/2$, $2L/3$ e $L/3$?

\paragraph{Modelo.}
\begin{equation}
  T = 2\pi\sqrt{\frac{L}{g}} \implies T \propto \sqrt{L}.
\end{equation}

\paragraph{Varia\c{c}\~ao relativa geral.}
\begin{equation}
  \frac{T'-T}{T} = \sqrt{\frac{L'}{L}}-1.
\end{equation}

\paragraph{Resultados.}
\begin{center}
\begin{tabular}{lll}
\toprule
Novo comprimento $L'$ & $T'/T$ & Varia\c{c}\~ao em $T$\\
\midrule
$L/2$    & $0{,}707$ & $-29{,}3\,\%$\\
$2L/3$   & $0{,}816$ & $-18{,}4\,\%$\\
$L/3$    & $0{,}577$ & $-42{,}3\,\%$\\
\bottomrule
\end{tabular}
\end{center}

% --- 5.4 ---
\subsection{Equa\c{c}\~oes de Lanchester -- Combate}

\subsubsection*{Modelo te\'orico (Lanchester, 1916)}

O sistema de EDOs acopladas descreve duas for\c{c}as $A(t)$ e $B(t)$:
\begin{equation}
  \frac{dA}{dt} = -\beta\,B, \qquad \frac{dB}{dt} = -\alpha\,A,
  \label{eq:lanchester}
\end{equation}
onde $\alpha$ e $\beta$ s\~ao as efic\'acias de combate (baixas por unidade por tempo).

\paragraph{Invariante de Lanchester.}
Multiplicando a 1\textordfeminine{} equa\c{c}\~ao por $2\alpha A$ e a 2\textordfeminine{}
por $2\beta B$ e somando:
\begin{equation}
  \frac{d}{dt}\!\left(\alpha\,A^2 - \beta\,B^2\right) = 0.
\end{equation}
Portanto o \textbf{invariante} \'e:
\begin{equation}
  \alpha\,A^2 - \beta\,B^2 = \alpha\,A_0^2 - \beta\,B_0^2 = K = \text{constante.}
  \label{eq:invariante}
\end{equation}

\paragraph{Crit\'erio de vit\'oria.}
\begin{itemize}
  \item $K > 0$: for\c{c}a $A$ vence; restam $A_f = \sqrt{K/\alpha}$ unidades.
  \item $K < 0$: for\c{c}a $B$ vence; restam $B_f = \sqrt{-K/\beta}$ unidades.
  \item $K = 0$: empate (ambas zeram simultaneamente).
\end{itemize}
\resultado{\alpha A_0^2 - \beta B_0^2 = K}

\subsubsection*{Exemplo: Batalha estil\'izada (dados hist\'oricos Iwo Jima)}

A Batalha de Iwo Jima (fev--mar 1945) \'e o caso cl\'assico de calibra\c{c}\~ao
das equa\c{c}\~oes de Lanchester (cf.\
Braun, \textit{Differential Equations and their Applications}, 4\textordfeminine{} ed.).

\begin{center}
\begin{tabular}{lll}
\toprule
For\c{c}a & Efetivo inicial & Efic\'acia de combate \\
\midrule
EUA\;($A$) & $A_0 = 54\,000$ & $\alpha = 0{,}0544\;\text{baixas/unidade/dia}$ \\
Jap\~ao\;($B$) & $B_0 = 21\,500$ & $\beta = 0{,}0106\;\text{baixas/unidade/dia}$ \\
\bottomrule
\end{tabular}
\end{center}

\paragraph{Passo 1: calcular o invariante.}
\begin{align}
  K &= \alpha\,A_0^2 - \beta\,B_0^2 \notag\\
    &= 0{,}0544\times(54\,000)^2 - 0{,}0106\times(21\,500)^2 \notag\\
    &= 0{,}0544\times 2{,}916\times10^{9} - 0{,}0106\times 4{,}622\times10^{8} \notag\\
    &= 1{,}586\times10^{8} - 4{,}899\times10^{6} \approx 1{,}537\times10^{8} > 0.
\end{align}
\resultado{K \approx 1{,}54\times10^{8} > 0 \Rightarrow \text{EUA vence}}

\paragraph{Passo 2: efetivo restante dos EUA ao final.}
Quando $B = 0$, da equa\c{c}\~ao \eqref{eq:invariante}:
\begin{equation}
  A_f = \sqrt{\frac{K}{\alpha}} = \sqrt{\frac{1{,}537\times10^{8}}{0{,}0544}}
      = \sqrt{2{,}825\times10^{9}} \approx 53\,153\;\text{soldados.}
\end{equation}
Baixas dos EUA: $54\,000 - 53\,153 \approx \mathbf{847}$ (pelo modelo puro de
atri\c{c}\~ao sim\'etrica --- na batalha real os EUA tiveram $\sim$26\,000 baixas
porque o Jap\~ao adotou t\'aticas de guerra de guerrilha, que quebram a hip\'otese
de combate direto do modelo).

\paragraph{Passo 3: quantos refor\c{c}os o Jap\~ao precisaria para empatar?}
No empate $K = 0$:
\begin{equation}
  \alpha\,A_0^2 = \beta\,B_{eq}^2
  \implies
  B_{eq} = A_0\sqrt{\frac{\alpha}{\beta}}
         = 54\,000\sqrt{\frac{0{,}0544}{0{,}0106}}
         = 54\,000\times 2{,}265 \approx 122\,310.
\end{equation}
O Jap\~ao precisaria de $\approx 122\,310 - 21\,500 = \mathbf{100\,810}$ refor\c{c}os
--- mais de $5\times$ seu efetivo inicial --- para empatar contra os EUA.

\paragraph{Interpreta\c{c}\~ao: lei quadr\'atica.}
Dobrar o efetivo multiplica o poder de combate por $2^2 = 4$ (n\~ao por $2$).
Isso explica a regra militar conhecida como \textbf{``concentra\c{c}\~ao de for\c{c}as''}:
\'e muito mais eficiente concentrar tropas do que distribu\'i-las.

\paragraph{Aplica\c{c}\~oes hist\'oricas confirmadas (Wikip\'edia, Lanchester's laws).}
\begin{itemize}
  \item \textit{Pickett's Charge} (Gettysburg, 1863): lei quadr\'atica prediz
    derrota Confederate --- confirmado historicamente.
  \item \textit{Battle of Britain} (1940): modelo calibrado com dados RAF vs.\ Luftwaffe.
  \item \textit{Battle of Kursk} (1943): melhor ajuste dos dados obtido com expoente $1{,}5$
    (entre linear e quadr\'atico), por causa de combate misto.
\end{itemize}

\subsubsection*{Se\c{c}\~ao 5.7 do Dym -- ``Dividir para Conquistar''? (An\'alise de Estrat\'egia)}

A Se\c{c}\~ao~5.7 do livro de Dym responde \`a pergunta:
\begin{quote}
  \itshape Compensa dividir as for\c{c}as e atacar o inimigo em etapas separadas,
  em vez de atacar de uma s\'o vez?
\end{quote}

\paragraph{Configura\c{c}\~ao do problema.}
\begin{itemize}
  \item For\c{c}a~$m$: exerce\'ito atacante com $m_0$ soldados.
  \item For\c{c}a~$n$: exer\'icio defensor com $n_0$ soldados; efic\'acias iguais
        ($\alpha = \beta = k$).
  \item Invariante (efic\'acias iguais): $m^2 - n^2 = \text{cte.}$
\end{itemize}

\paragraph{Estrat\'egia~I -- for\c{c}a $n$ unida.}
Ao final, quando $n = 0$:
\begin{equation}
  m_f^2 = m_0^2 - n_0^2 \implies m_f = \sqrt{m_0^2 - n_0^2}.
  \label{eq:unida}
\end{equation}

\paragraph{Estrat\'egia~II -- for\c{c}a $n$ dividida em dois grupos ($n_1 + n_2 = n_0$).}
\begin{enumerate}[label=\arabic*)]
  \item Batalha~1: $m_0$ vs.\ $n_1$:
    \begin{equation}
      m_1^2 = m_0^2 - n_1^2 \implies m_1 = \sqrt{m_0^2 - n_1^2}.
    \end{equation}
  \item Batalha~2 (sobreviventes $m_1$ vs.\ $n_2$):
    \begin{equation}
      m_2^2 = m_1^2 - n_2^2 = m_0^2 - n_1^2 - n_2^2.
      \label{eq:dividida}
    \end{equation}
\end{enumerate}

\paragraph{Compara\c{c}\~ao matem\'atica.}
Das equa\c{c}\~oes \eqref{eq:unida} e \eqref{eq:dividida}:
\begin{align}
  \text{Unida:}   \quad & m_f^2   = m_0^2 - (n_1+n_2)^2
                                   = m_0^2 - n_1^2 - 2n_1n_2 - n_2^2. \\[4pt]
  \text{Dividida:}\quad & m_2^2   = m_0^2 - n_1^2 - n_2^2.
\end{align}
Subtraindo:
\begin{equation}
  m_2^2 - m_f^2 = 2\,n_1\,n_2 > 0 \qquad (n_1, n_2 > 0).
\end{equation}
Logo $m_2 > m_f$: \textbf{mais soldados do exer\'icio $m$ sobrevivem quando o inimigo
$n$ divide suas for\c{c}as}.

\resultado{m_2^2 - m_f^2 = 2\,n_1\,n_2 > 0 \;\Rightarrow\; \text{Dividir \'e PIOR para }n}

\paragraph{Exemplo~1 -- $E_0 = 50\,000$,\; $F_0 = 70\,000$
  (for\c{c}a $F$ divide ao meio).}

A for\c{c}a $F$ tem superioridade num\'erica de $40\,\%$.
Pergunta: vale a pena $F$ atacar em dois grupos de $35\,000$ em vez de
atacar de uma s\'o vez?

\noindent\textbf{Cen\'ario~A -- $F$ unido.}
\begin{equation}
  E_f^2 = E_0^2 - F_0^2 = (50\,000)^2 - (70\,000)^2
        = 2{,}5\times10^{9} - 4{,}9\times10^{9} = -2{,}4\times10^{9} < 0.
\end{equation}
$E_f^2 < 0 \;\Rightarrow\;$ \textbf{$F$ vence};
sobreviventes de $F$:
$F_f = \sqrt{(70\,000)^2 - (50\,000)^2} = \sqrt{2{,}4\times10^{9}} \approx 48\,990$.

\noindent\textbf{Cen\'ario~B -- $F$ dividido em $F_1 = F_2 = 35\,000$.}

Usando a f\'ormula cumulativa $E_k^2 = E_0^2 - \displaystyle\sum_{i}F_i^2$:
\begin{align}
  \text{Batalha~1:}\quad
  E_1^2 &= (50\,000)^2 - (35\,000)^2 = 1{,}275\times10^{9}
  \;\implies\; E_1 \approx 35\,707, \\
  \text{Batalha~2:}\quad
  E_2^2 &= E_0^2 - F_1^2 - F_2^2
         = (50\,000)^2 - 2\times(35\,000)^2
         = 5{,}0\times10^{7}
  \;\implies\; E_2 \approx 7\,071.
\end{align}
$E_2^2 > 0 \;\Rightarrow\;$ \textbf{$E$ vence} com ${\approx}\,7\,071$ sobreviventes!

\begin{center}
\begin{tabular}{lcc}
\toprule
Estrat\'egia de $F$ & Vencedor & Sobreviventes finais \\
\midrule
$F$ unido ($70\,000$)              & $F$ & $F_f \approx 48\,990$ \\
$F$ dividido ($35\,000 + 35\,000$) & $E$ & $E_2 \approx 7\,071$ \\
\bottomrule
\end{tabular}
\end{center}

\textbf{Interpreta\c{c}\~ao}: dividir $F$ em dois grupos iguais
\emph{inverte completamente o resultado da batalha}.
A vantagem num\'erica de $40\,\%$ de $F$ \'e an\'ulada pelo efeito quadr\'atico
quando suas for\c{c}as s\~ao fragmentadas.

\paragraph{Exemplo~2 -- $E_0 = 90\,000$,\; $F_0 = 60\,000$
  (impacto de dividir em $k$ grupos iguais).}

Aqui $E$ j\'a \'e superior.
Analisa-se o que acontece \`a medida que $F$ divide seu efetivo em
$k$ grupos iguais de $60\,000/k$, cada um enfrentando $E$ sequencialmente.

\noindent\textbf{F\'ormula geral.}
Ap\'os $k$ batalhas sequenciais contra grupos $F_1, \ldots, F_k$
(com $\textstyle\sum_i F_i = F_0$):
\begin{equation}
  E_k^2 = E_0^2 - \sum_{i=1}^{k} F_i^2.
  \label{eq:geral_k}
\end{equation}
Para grupos iguais $F_i = F_0/k$:
\begin{equation}
  E_k^2 = E_0^2 - \frac{F_0^2}{k}.
\end{equation}

\begin{center}
\begin{tabular}{ccccc}
\toprule
$k$ & Tamanho do grupo & $E_k^2 = (90\,000)^2 - (60\,000)^2/k$
    & $E_k$ & Baixas de $E$ \\
\midrule
1        & $60\,000$ & $4{,}50\times10^{9}$ & $67\,082$ & $22\,918$ \\
2        & $30\,000$ & $6{,}30\times10^{9}$ & $79\,373$ & $10\,627$ \\
3        & $20\,000$ & $6{,}90\times10^{9}$ & $83\,066$ & $6\,934$  \\
6        & $10\,000$ & $7{,}50\times10^{9}$ & $86\,603$ & $3\,397$  \\
$\infty$ & $\to 0$   & $8{,}10\times10^{9}$ & $90\,000$ & $0$       \\
\bottomrule
\end{tabular}
\end{center}

\textbf{Interpreta\c{c}\~ao}: cada fragmenta\c{c}\~ao adicional reduz as baixas
infligidas a~$E$.
No limite $k \to \infty$ (ataques infinitesimais), $E$ n\~ao sofre baixa alguma
--- $F$ perde todo o seu poder de combate ao se dispersar.

\paragraph{Conclus\~ao (Dym, Se\c{c}\~ao 5.7).}
A lei quadr\'atica de Lanchester \emph{invalida} a m\'axima ``dividir para conquistar''.
A estrat\'egia \'otima \'e sempre \textbf{concentrar} as for\c{c}as.
Fragmentar o ataque desperdi\c{c}a o efeito quadr\'atico do tamanho, entregando ao
defensor uma vantagem crescente.

% --- 5.5 ---
\subsection{Juro Composto -- Problemas 5.9 e 5.47}

\paragraph{Problema 5.9 (Lista 4.34): dobrar em 7 anos.}
Para $(1+r)^7 = 2$:
\begin{equation}
  r = 2^{1/7}-1 = e^{\ln 2/7}-1 \approx e^{0{,}0990}-1 \approx 9{,}90\,\%\text{ ao ano.}
\end{equation}
\resultado{r \approx 9{,}90\,\%/\text{ano}}

\paragraph{Regra dos 72 (estimativa r\'apida).}
$r\approx 72/n = 72/7 \approx 10{,}3\,\%$ --- superestima em $\approx 4\,\%$ relativo.

\paragraph{Problema 5.47 (Lista 4.37): poupan\c{c}a de US\$\,1.000.000 a $5{,}5\%$\,a.a.}
\begin{equation}
  PV = \frac{FV}{(1+r)^n}.
\end{equation}
\begin{center}
\begin{tabular}{lll}
\toprule
Meta & Prazo & Dep\'osito em 2002\\
\midrule
US\$\,1\,000\,000 em 2022 & 20 anos & US\$\,342\,729\\
US\$\,1\,000\,000 em 2042 & 40 anos & US\$\,117\,463\\
\bottomrule
\end{tabular}
\end{center}
\resultado{PV_{40} \approx 0{,}34\,PV_{20}}

Dobrar o prazo reduz o dep\'osito necess\'ario em $\approx 66\,\%$ --- efeito direto da
exponen\-cia\c{c}\~ao do juro composto.

% ===================================================================
% RESUMO
% ===================================================================
\newpage
\section*{Resumo das respostas}
\addcontentsline{toc}{section}{Resumo das respostas}

\begin{center}
\begin{tabular}{lll}
\toprule
\textbf{Quest\~ao} & \textbf{Item} & \textbf{Resultado principal} \\
\midrule
Q1 & Conjunto~A (Koch)              & $D_f = \ln4/\ln3 \approx 1{,}26$ \\
Q1 & Conjunto~B (Sierpinski 3D)     & $D_f = \ln4/\ln2 = 2$ \\[4pt]
Q2 & $\Pi_1$                        & $L/h$ \\
Q2 & $\Pi_2$                        & $g/(h\omega^2)$ \\
Q2 & $\Pi_3$                        & $\tau/(mh\omega^2)$ \\
Q2 & Frequ\^encia                   & $\omega \propto \sqrt{g/h}$ e
                                       $\omega \propto \sqrt{\tau/(mh)}$ \\[4pt]
Q3a & Fun\c{c}\~ao                  & $V(t)=51{.}500\times(0{,}88)^t$ \\
Q3b & Taxa mensal                   & $i_m \approx 1{,}06\%/\text{m\^es}$ \\
Q3c & Gr\'afico                     & Exponencial decrescente; semi-log = reta \\
Q3d & $V(6\;\text{anos})$           & $\approx \text{R\$}\;23{.}917$ \\[4pt]
Q4a & Par m\'agico ($m=10$)         & $(a,c)=(31,9)$; condi\c{c}\~oes Hull-Dobell \\
Q4b & Demonstra\c{c}\~ao            & Per\'iodo $10$ verificado nas 3 sementes \\
\bottomrule
\end{tabular}
\end{center}

\end{document}
